% Capitulo: Dados como contexto para IA
\chapter{Dados como contexto para IA}
\label{cap:contexto-ia}

\rascunho{
  \textbf{Ideia do capítulo.} O uso mais comum de IA com dados nas empresas não é
  treinar, é dar contexto: o modelo recebe, na hora da pergunta, as definições, os
  documentos e os dados de que precisa para responder. A qualidade da resposta passa a
  depender de como esse contexto está organizado. É modelagem de dados com um
  consumidor novo.
  \begin{itemize}
    \item Ligação com \enquote{conversar com os dados} (\veja{sec:text-to-sql}) e com a
      camada semântica (\veja{sec:camada-semantica}).
    \item Ligação com custo: contexto grande custa tokens (ver o
      capítulo~\ref{cap:custo-consumo}).
  \end{itemize}
}

% ---------------------------------------------------------------------------
\section{O modelo só sabe o que recebe}
\label{sec:modelo-so-sabe}

\rascunho{
  \begin{itemize}
    \item Um modelo de linguagem não conhece as tabelas, as métricas nem os documentos da
      empresa. Ele responde com o que está no contexto, e inventa quando falta.
    \item Duas formas principais de dar contexto: buscar trechos de documentos
      relevantes para a pergunta (geração aumentada por recuperação, RAG)
      \parencite{lewis2020}, e dar acesso a dados estruturados com as definições de
      negócio junto.
  \end{itemize}
}

\pendente{escrever}

% ---------------------------------------------------------------------------
\section{Metadados e camada semântica como contexto}
\label{sec:metadados-contexto}

\rascunho{
  \begin{itemize}
    \item Descrição de cada tabela e coluna, definição de cada métrica, relacionamentos,
      exemplos de consultas certas: tudo que o analista novo precisaria saber, a IA
      também precisa.
    \item A documentação que antes era \enquote{boa prática} vira o insumo principal.
      Mesmo raciocínio do TDD (\veja{sec:tdd}): a IA tornou barato e necessário o que
      antes ficava para depois.
    \item Um modelo dimensional bem feito, com nomes claros, já é meio caminho: grão
      explícito e métricas definidas uma vez reduzem a chance de a IA somar errado.
  \end{itemize}
}

\pendente{escrever}

% ---------------------------------------------------------------------------
\section{Documentos e recuperação}
\label{sec:documentos-recuperacao}

\rascunho{
  \begin{itemize}
    \item Documentos divididos em trechos, com metadados (origem, data, dono, nível de
      acesso), e um índice de busca por significado (vetores, \emph{embeddings}).
    \item O mesmo raciocínio do medalhão: o documento bruto, o documento limpo e
      dividido, o índice pronto para consulta (ver a Figura~\ref{fig:medalhao}).
    \item Permissões: a IA não pode responder com um documento que quem perguntou não
      poderia ler.
    \item Atualização: contexto velho dá resposta velha com cara de nova. Vale o
      capítulo~\ref{cap:frequencia} inteiro.
  \end{itemize}
}

\pendente{escrever}

% ---------------------------------------------------------------------------
\section{Organizar para a IA}
\label{sec:organizar-para-ia}

\rascunho{
  \begin{itemize}
    \item Uma lista prática: catálogo com dono para cada fonte, glossário de negócio,
      métricas definidas uma vez, exemplos de perguntas com a resposta certa, permissões
      herdadas da fonte.
    \item Avaliar as respostas com um conjunto de perguntas conhecidas (ver
      \veja{sec:consumidor-modelo}), a cada mudança de dado, de contexto ou de modelo.
    \item Relato: o que eu precisei organizar para uma IA responder perguntas sobre os
      nossos dados. \pendente{escrever com relato, sem nomes}
  \end{itemize}
}

\pendente{escrever}

% ---------------------------------------------------------------------------
\begin{leituras}
  \leitura{lewis2020}{O artigo que propôs a geração aumentada por recuperação (RAG),
    base da maior parte dos assistentes que respondem com documentos da empresa.}
  \leitura{huyen2022}{Para a visão de sistema: como dados, modelo e avaliação se
    encaixam em produção.}
\end{leituras}
